\begin{abstract}
跨越互不完全信任组织的数据流通场景，不只需要文件级完整性校验。实际可部署的审计系统还需要支持可追踪授权、可验证链下证据、可问责恢复，以及高频审计事件下的低成本确认。BLoP 提供了一种融合区块链、可信计算、隐私保护计算、异常监测和 PBFT 风格确认机制的全节点可信审计模型。本文研究该模型的一种轻量化扩展 TrustAuditFlow，面向审计事件频繁且并非所有事件都需要全节点确认的数据流通场景。TrustAuditFlow 保留 BLoP 作为强安全基线，同时增加风险感知的委员会路径：低风险和中风险审计事件在链下聚合，通过 Merkle 根和哈希链锚定，并由可信审计委员会确认；高风险或存在争议的事件则回退到 BLoP 风格的全节点路径。该设计进一步将数据分片、抽样完整性校验、基于纠删码的恢复、链上问责记录和信任分更新连接为闭环审计恢复流程，为数据流通中的区块链分布式数据校验提供结构化工作流。
\end{abstract}

\begin{IEEEkeywords}
数据流通，区块链审计，分布式数据校验，可信委员会，PDP，PoR，纠删码，PBFT，Merkle 树
\end{IEEEkeywords}

